\documentclass[12pt,a4paper]{article}
\usepackage[utf8]{inputenc}
\usepackage{amsmath}
\usepackage{amssymb}
\usepackage{amsthm}
\usepackage{geometry}
\usepackage{hyperref}
\usepackage{graphicx}

\geometry{margin=1in}

\title{\textbf{Commons Production}\\
\large A System for Transparent Labor-Time Accounting}
\author{}
\date{}

\begin{document}

\maketitle

\begin{abstract}
We present a complete system for communal production based on transparent labor-time accounting. Beginning from the foundational premise that the social character of production is presupposed, we develop the theoretical framework, implement transparent accounting through the Stock-Book, and show how planned allocation realizes the first economic law. Finally, we demonstrate that distribution of means of consumption is a necessary consequence of the organization of production itself.
\end{abstract}

\section{The Foundational Premise}

\subsection{Social Character Presupposed}

In communal production, the social character of labor is \textit{presupposed} from the outset—not achieved through exchange.\footnote{``Within the co-operative society based on common ownership of the means of production, the producers do not exchange their products; just as little does the labor employed on the products appear here as the \textit{value} of these products.'' --- Marx, \textit{Critique of the Gotha Programme}}

\begin{equation}
\text{Individual Labor} = \text{Social Labor (from outset)}
\end{equation}

Products are communal and general from their creation. They never become commodities; they never acquire exchange value.
\subsection{Total Social Product}

The combined labor power of the community, applied consciously, produces the total social product:\footnote{``Before division among the individuals, there has to be deducted from it: First, cover for replacement of the means of production used up. Second, additional portion for expansion of production. Third, reserve or insurance funds...'' --- Marx, \textit{Critique of the Gotha Programme}}

\begin{equation}
\text{Total Social Product} = \sum(\text{All Products of Labor})
\end{equation}

This must be decomposed before individual distribution.

\subsection{The Six Deductions}

\subsubsection{Production-Related (Economic Necessity)}

\begin{align}
D_1 &= \text{Replacement of means of production used up} \\
D_2 &= \text{Additional portion for expansion of production} \\
D_3 &= \text{Reserve/insurance funds (accidents, calamities)}
\end{align}

These are determined by available means and probability calculations, \textit{not by equity}:

\begin{equation}
\text{Magnitude}(D_1, D_2, D_3) = f(\text{Available Means}, \text{Probabilities})
\end{equation}

\subsubsection{Social/Administrative (Development-Dependent)}

From the remainder—the means of consumption—we deduct:\footnote{``costs of administration not belonging to production... that which is intended to serve the common satisfaction of needs, such as schools, health services, etc.'' --- Marx, \textit{Critique of the Gotha Programme}}

\begin{align}
D_4 &= \text{General administration (not production)} \\
D_5 &= \text{Common needs (schools, health, etc.)} \\
D_6 &= \text{Funds for those unable to work}
\end{align}

These exhibit reciprocal dynamics with social development:

\begin{align}
\Delta(D_4) &\propto -\Delta(\text{Development}) \quad \text{(admin shrinks)} \\
\Delta(D_5) &\propto +\Delta(\text{Development}) \quad \text{(common needs grow)}
\end{align}

\subsubsection{Distribution Pool}

\begin{equation}
\text{Distribution Pool} = \text{Total Product} - \sum_{i=1}^{6} D_i
\end{equation}

\section{Economy of Time: The First Economic Law}

\subsection{The Fundamental Principle}

All economy resolves to economy of time:\footnote{``Economy of time, to this all economy ultimately reduces itself.'' --- Marx, \textit{Grundrisse}}

\begin{equation}
\boxed{\text{All Economy} = \text{Economy of Time}}
\end{equation}

Wealth is proportional to free time:\footnote{``Real economy — saving — consists of the saving of labour time.'' --- Marx, \textit{Grundrisse}}

\begin{equation}
T_{\text{free}} = T_{\text{total}} - T_{\text{necessary}}
\end{equation}

More wealth means less necessary labor time:

\begin{equation}
\Delta(\text{Wealth}) \propto -\Delta(T_{\text{necessary}})
\end{equation}

\subsection{Distributing Time Across Branches}

Society must consciously distribute its time across different kinds of work:\footnote{``Economy of time and planned distribution of labour time among the various branches of production... [this] becomes the first economic law on the basis of communal production.'' --- Marx, \textit{Grundrisse}}

\begin{equation}
\sum_{i} t_i = T_{\text{total social time}}
\end{equation}

Where allocation depends on difficulty, importance, and available time:

\begin{equation}
t_i = f(\text{difficulty}_i, \text{importance}_i, T_{\text{available}})
\end{equation}

\section{The Stock-Book: Making Production Transparent}

\subsection{Simple and Intelligible Relations}

The Stock-Book makes social relations transparent:\footnote{``The social relations of the individual producers, with regard both to their labor and to its products, are in production as well as in distribution \textit{simple and intelligible}.'' --- Marx, \textit{Capital, Vol. I}}

\begin{equation}
\text{Stock-Book Records} = \{\text{Products, Operations, Average Labor-Time}\}
\end{equation}

Every operation records inputs (products and labor) and effects (state changes):

\begin{table}[h]
\centering
\small
\begin{tabular}{lllll}
\hline
Op ID & Total Time & Inputs (Products) & Inputs (Labor) & Effects \\
\hline
OP056 & 8h & Wood:10kg, Saw:8h & I001(6h), I002(2h) & Chairs:+2 \\
OP057 & 4h & Classroom:4h & I003(4h) & LP\_engineer:+20h \\
\hline
\end{tabular}
\end{table}

\subsection{Average Labor-Time (ALT)}

Every product has an average labor-time embodied in it:

\begin{equation}
\text{ALT} = \frac{\sum(\text{Total Labor-Time})}{\sum(\text{Quantity Produced})}
\end{equation}

\subsection{Using-Up: The Universal Formula}

Labor-time transfers from products to operations:

\begin{equation}
LT_{\text{transferred}} = \frac{\text{Usage}}{\text{Lifespan}(P)} \times \text{ALT}(P)
\end{equation}

\subsection{Labor Power: The Dual System}

Labor power requires two orthogonal accounting systems:

\textbf{System 1: Daily Energy} - Time scale: Daily (24h cycles)

\textbf{System 2: Skill Depreciation} - Time scale: Years/Decades

These systems are mathematically independent.

\subsection{Multi-Product Allocation}

When operations produce multiple products:

\begin{align}
\text{Total Input LT} &= \sum(\text{Product Inputs}) + \sum(\text{Labor Inputs}) \\
\text{ALT}_A &= \frac{a_A \times \text{Total Input LT}}{\text{qty}_A}
\end{align}

Allocation weights $a_A$ are social choices.

\subsection{Rolling Averages}

New production updates existing ALT:

\begin{equation}
\text{New ALT}(P) = \frac{Q_{\text{old}} \times \text{ALT}_{\text{old}} + LT_{\text{allocated}}}{Q_{\text{old}} + q_{\text{new}}}
\end{equation}

\subsection{Consumption as Production}

Every consumption operation produces restored labor-power:

\begin{equation}
\text{Products} + \text{Time} \rightarrow \text{Restored LP} \rightarrow \text{Labor} \rightarrow \text{Products}
\end{equation}

\subsection{Conservation}

\begin{equation}
\sum(\text{Input LT}) = \sum(\text{Output LT}) + \text{Waste LT}
\end{equation}

Every hour of social labor is traceable.

\section{The Allocation System: Planning Production}

\subsection{Allocation Records}

The Stock-Book records reality; the Allocation system plans the future:

\begin{table}[h]
\centering
\small
\begin{tabular}{llllll}
\hline
Period & Product & Planned & Actual & Variance & Status \\
\hline
2024-W10 & P001 & 100kg & 98kg & -2\% & Completed \\
2024-W10 & LP\_base & 40h & 42h & +5\% & Overused \\
\hline
\end{tabular}
\end{table}

\subsection{Buffers and Reserves}

\begin{equation}
\text{Planned Allocation} = \text{Expected Need} \times (1 + \text{buffer})
\end{equation}

\subsection{The Planning Hierarchy}

\begin{enumerate}
\item \textbf{Annual}: Set deduction targets ($D_1$ through $D_6$)
\item \textbf{Monthly}: Allocate products to operations with buffers
\item \textbf{Weekly}: Execute, record in Stock-Book, track variance
\end{enumerate}

\subsection{The Complete Planning Problem}

\textbf{Objective}: Maximize free time

\begin{equation}
\text{Maximize:} \quad \sum_t \left( T_{\text{total},t} - T_{\text{necessary},t} \right)
\end{equation}

\textbf{Constraints}: Stock balance, deduction requirements, labor balance

\textbf{Result}: Good planning reduces variance and shortages, increasing free time.

\subsection{Dynamic Learning}

\begin{equation}
\text{Plan} \rightarrow \text{Execute} \rightarrow \text{Record} \rightarrow \text{Analyze} \rightarrow \text{Adjust} \rightarrow \text{Replan}
\end{equation}

When variance exceeds threshold:

\begin{align}
\text{Expected Need}_{\text{new}} &= \text{Rolling Avg}(\text{Actual}) \\
\text{Buffer}_{\text{new}} &= f(\text{New Variance})
\end{align}

\section{Distribution as Consequence of Production Relations}

\subsection{The General Principle}

Mode of distribution varies based on productive organization and historical development:\footnote{``The mode of distribution of the means of consumption is only a consequence of the distribution of the conditions of production themselves. The latter distribution, however, is a feature of the mode of production itself.'' --- Marx, \textit{Critique of the Gotha Programme}}

\begin{equation}
\text{Mode of Distribution} = f(\text{Productive Organization}, \text{Historical Development})
\end{equation}

\subsection{The Fundamental Relation}

\begin{align}
\text{Distribution of Consumption} &= \text{Consequence of Distribution of Production Conditions} \\
\text{Distribution of Production Conditions} &= \text{Feature of Mode of Production}
\end{align}

Thus:

\begin{equation}
\text{Mode of Production} \rightarrow \text{Distribution of Conditions} \rightarrow \text{Distribution of Consumption}
\end{equation}

\textbf{Capitalist Mode}: Material conditions in hands of non-workers (capital + land) $\rightarrow$ masses own only labor power $\rightarrow$ present-day consumption distribution results automatically.

\textbf{Co-operative Mode}: Material conditions = co-operative property of workers $\rightarrow$ different consumption distribution results.

\begin{equation}
\text{Co-operative Ownership} \rightarrow \text{Different Distribution}
\end{equation}

\subsection{Labor Time: The Dual Role}

Labor time plays two distinct roles:

\subsubsection{First Role: Planning and Apportionment}

Labor time apportionment in accordance with a definite social plan maintains proper proportion between different kinds of work and the various wants of the community:

\begin{equation}
\text{Apportionment ensures: } \sum(\text{labor time}_i) \text{ achieves production adequate for } \sum(\text{want}_i)
\end{equation}

\subsubsection{Second Role: Measuring Contribution}

Labor time measures individual contribution:

\begin{align}
\text{Individual's Portion} &= \frac{\text{Individual Labor Time}}{\text{Total Labor Time}} \\
\text{Individual's Share of Individual Consumption Pool} &\propto \text{Individual's Labor Time}
\end{align}

\subsection{Labor Certificates and Equal Exchange}

\subsubsection{The Exchange Principle}

Individual producer receives back from society exactly what is given (after deductions):\footnote{``The individual producer receives back from society — after the deductions have been made — exactly what he gives to it... He receives a certificate from society that he has furnished such-and-such an amount of labor.'' --- Marx, \textit{Critique of the Gotha Programme}}

\begin{align}
\text{Social Working Day} &= \sum_i \text{Individual Hours}_i \\
\text{Individual's Share} &= \frac{\text{Individual Labor Time}}{\text{Social Working Day}}
\end{align}

\subsubsection{Labor Certificates Mechanism}

Individual receives certificate showing labor furnished (after deductions for common funds). With this certificate, draws from social stock of means of consumption an amount equal to labor cost.

\begin{equation}
\text{Labor Given (one form)} = \text{Labor Received (another form)}
\end{equation}

\subsubsection{Principle of Equal Values}

Same principle as exchange of commodities (exchange of equal values), but content and form changed:

\begin{itemize}
\item No one can give anything except labor
\item Nothing can pass to individuals except individual means of consumption
\item Given amount of labor in one form = exchanged for equal amount in another form
\end{itemize}

\textbf{Equal right in principle = bourgeois right}, but principle and practice no longer at loggerheads. Exchange of equivalents exists in individual case (not just on average).

\subsection{Measurement of Labor: Standards and Inequality}

\subsubsection{The Equal Standard}

Right of producers is proportional to labor supplied:\footnote{``The right of the producers is \textit{proportional} to the labor they supply; the equality consists in the fact that measurement is made with an \textit{equal standard}, labor.'' --- Marx, \textit{Critique of the Gotha Programme}}

\begin{equation}
\text{Equality} = \text{Measurement with Equal Standard (Labor)}
\end{equation}

Labor as measure must be defined by duration or intensity. Without this definition, labor ceases to be a standard.

\subsubsection{Qualitative vs Quantitative Differences}

Labor differs quantitatively \textit{and} qualitatively. Different kinds of labor are direct social functions (immediately social) within a spontaneously developed division of labor.

Quantitative comparison requires qualitative equivalence. Measuring labor quantitatively assumes the identity of their quality.

\subsubsection{Unequal Individuals, Equal Standard}

Equal right = unequal right for unequal labor:\footnote{``But one man is superior to another physically, or mentally, and supplies more labor in the same time, or can labor for a longer time... This \textit{equal} right is an \textit{unequal} right for \textit{unequal} labor.'' --- Marx, \textit{Critique of the Gotha Programme}}

\begin{itemize}
\item Recognizes no class differences (everyone is only a worker)
\item Tacitly recognizes unequal individual endowment + productive capacity as natural privilege
\item Right = right of inequality (in its content)
\item Unequal individuals measurable by equal standard only when brought under equal point of view
\item Individuals regarded only as workers (everything else ignored)
\end{itemize}

\subsubsection{Inequality in Outcomes}

Equal performance of labor + equal share in social consumption fund $\rightarrow$ unequal actual receipt.

One worker married, another not; one has more children, etc. With equal labor performance: one receives more than another, one richer than another.

To avoid these defects, right would have to be unequal (instead of equal).

\subsubsection{Constraints of First Phase}

These defects are inevitable in the first phase of communist society:\footnote{``Right can never be higher than the economic structure of society and its cultural development conditioned thereby.'' --- Marx, \textit{Critique of the Gotha Programme}}

\begin{equation}
\text{Right} \leq f(\text{Economic Structure}, \text{Cultural Development})
\end{equation}

\subsection{Higher Phase of Communist Society}

\subsubsection{Conditions for Higher Phase}

\begin{itemize}
\item Enslaving subordination to division of labor has vanished
\item Antithesis between mental and physical labor has vanished
\item Labor has become: not only means of life, but life's prime want
\item Productive forces increased with all-around development of individual
\item Springs of co-operative wealth flow more abundantly
\end{itemize}

\subsubsection{The Principle}

Then the narrow horizon of bourgeois right is crossed in its entirety and society may inscribe on its banner:

\begin{equation}
\boxed{\text{From each according to ability, to each according to needs}}
\end{equation}

\section{Conclusion}

We have presented a complete system that realizes Marx's vision:

\begin{enumerate}
\item The social character of production is \textit{presupposed}
\item Economy of time plus planned distribution = the first economic law
\item The Stock-Book makes all labor-time transparent and traceable
\item Planned allocation organizes production to maximize free time
\item Distribution of consumption is a \textit{consequence} of organizing production co-operatively
\end{enumerate}

The key insight: once we organize the \textit{conditions of production} co-operatively (Stock-Book + Allocation), the distribution of \textit{means of consumption} follows necessarily. Distribution is not an independent problem to be solved—it is the automatic result of how production itself is organized.

In the first phase, labor time plays its dual role: planning production and measuring contribution. In the higher phase, when productive forces have developed sufficiently, society can transcend bourgeois right entirely and realize: from each according to ability, to each according to needs.

\end{document}
